% chapters/ch21-conclusion-next-steps.tex
\chapter{Conclusion and Next Steps}

\begin{learningobjectives}
In this chapter you will consolidate the conceptual design workflow.
You will see what is required to move toward implementation later.
You will identify the artifacts needed for a clean handoff.
\end{learningobjectives}

\section{What you have built}
You now have a conceptual model grounded in meaning.
It includes entities, relationships, and constraints.
It is independent of any SQL dialect.

\section{From model to implementation}
Implementation is a translation step, not a redesign.
The model provides the rules.
The implementation provides the mechanics.

\subsection{Mapping as a disciplined translation}
Entities become relations.
Relationships become references or association relations.
Constraints become enforceable checks.

\section{Documentation as a deliverable}
A model is only useful if it is shared and understood.
Your documentation should be precise and minimal.

\subsection{Core artifacts}
At minimum, provide:
\begin{itemize}
  \item A diagram for communication.
  \item A rule list with explicit constraints.
  \item A glossary of terms and relationships.
\end{itemize}

\begin{examplebox}
\textbf{Handoff package}\
Diagram, rule list, constraint catalog, and a short glossary.\\
These are enough for implementation without ambiguity.
\end{examplebox}

\section{Readiness checks}
Before implementation, confirm these points.
They prevent late surprises.

\subsection{Model stability}
Do key entities and rules remain consistent under edge cases?
If not, return to requirements.

\subsection{Stakeholder agreement}
A model is valid only when stakeholders agree with the rules.
Review is part of design, not an afterthought.

\begin{principlebox}
\textbf{Principle}\
Clarity comes before optimization.
\end{principlebox}

\section{What comes next (but not yet)}
Implementation will introduce SQL, indexes, and performance choices.
Those topics are important, but they belong to a later phase.
For now, focus on meaning and correctness.

\begin{examplebox}
\textbf{Tease the next phase}\
You will later choose keys, constraints, and indexes in SQL.\\
Those choices should mirror the conceptual rules, not replace them.
\end{examplebox}

\section{Closing loop}
A good model is a living artifact.
It should be revisited as the domain evolves.
Treat it as part of the system, not a one-off document.

\begin{chaptersummary}
You have a complete conceptual framework for relational design.
The next step is implementation, but only after clarity and agreement.
Documentation, rules, and constraints form a reliable handoff.
\end{chaptersummary}

\begin{exercisebox}
Reflect on a model you have built or reviewed.
\begin{enumerate}[label=\arabic*.]
  \item List three rules that were most critical to meaning.
  \item Identify one place where the model was still ambiguous.
  \item Write a short handoff checklist for a future implementation team.
  \item Describe one risk if implementation starts too early.
\end{enumerate}
\end{exercisebox}
